\chapter{Ergebnisse}
\label{chap:ergebnisse}

\section{F1 — Prozentuale Zeitreduktion}

Tabelle~\ref{tab:q1} zeigt die prozentuale Zeitreduktion aller 17~Spiele, sortiert nach Gesamtreduktion absteigend.

\begin{table}[htb]
  \centering
  \caption{Prozentuale Zeitreduktion und jährliche Verbesserungsrate je Spiel (any\%-Kategorie)}
  \label{tab:q1}
  \small
  \begin{tabular}{llrr}
    \toprule
    \textbf{Spiel} & \textbf{Genre} & \textbf{Reduktion (\%)} & \textbf{Jährl.\ Rate (\%/a)} \\
    \midrule
    Minecraft: Java Edition         & Sandbox         & 99,5 & 8,9  \\
    The Talos Principle             & Puzzle          & 95,9 & 12,8 \\
    Portal                          & Puzzle          & 91,1 & 6,7  \\
    Hollow Knight                   & Action-Adventure& 89,8 & 12,8 \\
    Celeste                         & Platformer      & 80,2 & 9,8  \\
    Half-Life~2                     & FPS             & 75,4 & 6,2  \\
    Zelda: Ocarina of Time          & Action-Adventure& 71,2 & 15,8 \\
    Pac-Man                         & Arcade          & 69,0 & 11,2 \\
    Doom~3                          & FPS             & 61,6 & 4,8  \\
    Portal~2                        & Puzzle          & 55,2 & 4,0  \\
    Quake                           & FPS             & 47,5 & 1,8  \\
    Super Metroid                   & Action-Adventure& 41,5 & 3,1  \\
    Final Fantasy~VII               & RPG             & 34,3 & 2,9  \\
    Super Mario~64                  & Platformer      & 27,1 & 1,4  \\
    Donkey Konga                    & Arcade          & 14,9 & 2,0  \\
    Pokémon Red/Blue                & RPG             & 14,0 & 1,0  \\
    Super Mario Bros.               & Platformer      & 9,4  & 0,4  \\
    \bottomrule
  \end{tabular}
\end{table}

Auf Genre-Ebene zeigen sich deutliche Unterschiede in der mittleren Gesamtreduktion: \textit{Sandbox}~(99,5~\%), \textit{Puzzle}~(80,7~\%), \textit{Action-Adventure}~(67,5~\%) und \textit{FPS}~(61,5~\%) liegen deutlich über \textit{Arcade}~(41,9~\%), \textit{Platformer}~(38,9~\%) und \textit{RPG}~(24,2~\%). Der Kruskal-Wallis-Test auf die \emph{jährlichen Verbesserungsraten} ergibt jedoch keinen signifikanten Unterschied zwischen den sieben~Genres ($H = 6{,}24$; $p = 0{,}284$).

\finding{1}{Die prozentuale Zeitreduktion variiert stark zwischen Spielen (\textbf{9,4~\% bis 99,5~\%}). Spiele mit kurzer Spielzeit oder vielen nutzbaren Programmfehlern (\textit{Minecraft}, Puzzle-Spiele) erzielen die höchsten Reduktionen.}

\finding{2}{Ein Kruskal-Wallis-Test auf die jährlichen Verbesserungsraten zeigt keinen statistisch signifikanten Unterschied zwischen Genres ($\mathbf{H = 6{,}24}$; $\mathbf{p = 0{,}284}$). Die Verbesserungsrate hängt demnach stärker vom einzelnen Spiel als vom Genre ab.}

\section{F2 — Sättigung der Verbesserungsrate}

Tabelle~\ref{tab:q2_modelle} fasst die Modellvergleiche zusammen. Über alle 17~Spiele gewinnt das quadratische Polynom~(poly2) am häufigsten nach AIC~(6/17~Spiele), gefolgt von Gompertz~(5/17), exponentiellem Zerfall~(4/17) und logarithmischer Regression~(2/17). Der mittlere $R^2$-Wert des jeweils besten Modells beträgt~$0{,}93$ (Bereich: $0{,}73$–$0{,}98$).

\begin{table}[htb]
  \centering
  \caption{AIC-Modellvergleich: Siege und mittleres $R^2$ je Modelltyp (alle 17~Spiele)}
  \label{tab:q2_modelle}
  \small
  \begin{tabular}{lrr}
    \toprule
    \textbf{Modell} & \textbf{Siege (von 17)} & \textbf{Ø $R^2$} \\
    \midrule
    Quadr.\ Polynom (poly2)  & 6 & 0,901 \\
    Gompertz                  & 4 & 0,961 \\
    Exponentieller Zerfall    & 4 & 0,950 \\
    Logarithmisch             & 2 & 0,929 \\
    Potenzgesetz (power\_law) & 1 & 0,720 \\
    \bottomrule
  \end{tabular}
\end{table}

Der Chow-Test zeigt für \emph{alle 17~Spiele} einen statistisch signifikanten Strukturbruch. Tabelle~\ref{tab:chow} listet Bruchpunkt und $F$-Statistik je Spiel; alle $F$-Werte liegen deutlich über dem kritischen Wert ($p < 0{,}001$). Bei 8~von 17~Spielen liegt der Bruchpunkt innerhalb der ersten 10~Weltrekorde, also sehr früh in der Community-Entwicklung.

\begin{table}[htb]
  \centering
  \caption{Chow-Test: Bruchpunkt (WR-Nummer) und $F$-Statistik je Spiel (alle $p < 0{,}001$)}
  \label{tab:chow}
  \small
  \begin{tabular}{llrr}
    \toprule
    \textbf{Spiel} & \textbf{Genre} & \textbf{Bruchpunkt} & \textbf{$F$} \\
    \midrule
    Celeste                         & Platformer       & WR\,\#6  &    105,5 \\
    Donkey Konga                    & Arcade           & WR\,\#8  &     52,1 \\
    Doom~3                          & FPS              & WR\,\#24 &    881,3 \\
    Final Fantasy~VII               & RPG              & WR\,\#28 &     56,5 \\
    Half-Life~2                     & FPS              & WR\,\#50 &     66,4 \\
    Hollow Knight                   & Action-Adventure & WR\,\#33 &     33,6 \\
    Minecraft: Java Edition         & Sandbox          & WR\,\#7  &     99,3 \\
    Pac-Man                         & Arcade           & WR\,\#4  &     54,7 \\
    Pokémon Red/Blue                & RPG              & WR\,\#20 &     54,1 \\
    Portal                          & Puzzle           & WR\,\#4  & 20162,0  \\
    Portal~2                        & Puzzle           & WR\,\#28 &     64,3 \\
    Quake                           & FPS              & WR\,\#5  &     90,9 \\
    Super Mario Bros.               & Platformer       & WR\,\#10 &    410,0 \\
    Super Mario~64                  & Platformer       & WR\,\#16 &    118,9 \\
    Super Metroid                   & Action-Adventure & WR\,\#53 &     20,4 \\
    The Talos Principle             & Puzzle           & WR\,\#23 &    161,8 \\
    Zelda: Ocarina of Time          & Action-Adventure & WR\,\#10 &    488,4 \\
    \bottomrule
  \end{tabular}
\end{table}

\finding{3}{Für \textbf{alle 17~Spiele} gelingt eine Kurvenanpassung mit $\mathbf{R^2 > 0{,}70}$ (Median: $0{,}95$). Kein einzelnes Modell dominiert universell; \textit{poly2}~(6/17) erzielen die meisten Siege nach AIC, gefolgt von \textit{Gompertz} und exponentiellem Zerfall (je 4/17). Modelle mit explizitem Sättigungsverhalten (Gompertz, Exp.\ Zerfall) gewinnen gemeinsam für 8/17~Spiele.}

\finding{4}{\textbf{Alle 17~Spiele} weisen einen statistisch signifikanten Strukturbruch auf (Chow-Test, $p < 0{,}05$). Bei \textbf{8 von 17~Spielen} tritt der Bruchpunkt innerhalb der ersten 10~Weltrekorde auf; bei Spielen mit langer WR-Geschichte liegt er deutlich später (Median: WR~\#16). In beiden Fällen folgt auf den initialen Optimierungsschub (Routing, \emph{Glitch}-Entdeckungen) eine flachere \emph{Sättigungsphase}.}

\section{F3 — Externer Einfluss auf die Rekordentwicklung}

\subsection*{F3a — Post-Breakthrough-Dynamik}

Tabelle~\ref{tab:f3a} zeigt das Flattening Ratio sowie Breakthrough-Kennzahlen für alle Spiele, bei denen Pre- und Post-Breakthrough-Segmente ermittelt werden konnten.

\begin{table}[htb]
  \centering
  \caption{Post-Breakthrough-Dynamik: Flattening Ratio $\phi$ und Post-Trend~$\rho$ je Spiel (any\%-Kategorie, $n = 6$)}
  \label{tab:f3a}
  \small
  \begin{tabular}{lrrrll}
    \toprule
    \textbf{Spiel} & \textbf{Break (s)} & \textbf{\% Ges.} & \textbf{$\phi$} & \textbf{$\rho$} & \textbf{Interpretation} \\
    \midrule
    Donkey Konga                 & 291,0 & 42,4 & 23,97 & n.\,a.               & Beschleunigung \\
    Pokémon Red/Blue             & 138,0 & 13,6 &  0,44 & $+$0,31              & Plateau \\
    Portal~2                     & 432,8 & 16,1 &  0,35 & $-$0,54$^{*}$        & Plateau \\
    Quake                        & 183,0 & 29,4 &  0,04 & $-$0,49              & Starkes Plateau \\
    Zelda: Ocarina of Time       & 116,2 & 20,6 &  0,02 & $-$0,79$^{***}$      & Starkes Plateau \\
    The Talos Principle          & 864,0 & 19,0 &  0,02 & $-$0,79$^{***}$      & Starkes Plateau \\
    \bottomrule
  \end{tabular}
  \smallskip\par\noindent
  {\small \% Ges.\ = Breakthrough-Anteil an der Gesamtreduktion. $^{*}$\,$p < 0{,}05$; $^{***}$\,$p < 0{,}001$; n.\,a.\ = zu wenige Post-WRs für Spearman-Test. 11 von 17~Spielen hatten unzureichende Segmentgröße für die Velocity-Berechnung.}
\end{table}

\finding{5}{\textbf{5 von 6} auswertbaren Spielen zeigen ein Plateau nach dem Breakthrough ($\phi < 0{,}5$). Nur \textit{Donkey Konga} zeigt eine Beschleunigung ($\phi = 23{,}97$) --- der Breakthrough (WR~\#9, Dezember~2024) liegt so kurz zurück, dass die Post-Phase nur wenige Folgeverbesserungen enthält. Die stärksten Plateaus weisen \textit{Zelda: Ocarina of Time} und \textit{The Talos Principle} auf (je $\phi = 0{,}02$; $\rho = -0{,}79$; $p < 0{,}001$). Der Breakthrough-Anteil an der Gesamtreduktion reicht von \textbf{13,6~\%} (\textit{Pokémon Red/Blue}) bis \textbf{42,4~\%} (\textit{Donkey Konga}).}

\subsection*{F3b — Annäherung ans theoretische Minimum}

Tabelle~\ref{tab:f3b_proxy} zeigt den modellbasierten Asymptoten-Proxy für alle 17~Spiele.

\begin{table}[htb]
  \centering
  \caption{Modellbasierter Asymptoten-Proxy: Theoretisches Minimum aus exp\_decay-Modell ($R^2 \geq 0{,}70$)}
  \label{tab:f3b_proxy}
  \small
  \begin{tabular}{lrrr}
    \toprule
    \textbf{Spiel} & \textbf{Modell-Floor (s)} & \textbf{Gap (s)} & \textbf{\% Reduktion erreicht} \\
    \midrule
    Half-Life~2         & 1874,7 & 315,6 & 95,5 \\
    Pokémon Red/Blue    & 6004,4 & 215,6 & 82,5 \\
    \midrule
    \multicolumn{4}{l}{\textit{Weitere 15 Spiele: \texttt{none\_detected} --- kein konvergentes Sättigungsmodell ($R^2 < 0{,}70$)}} \\
    \bottomrule
  \end{tabular}
\end{table}

\finding{6}{Für \textbf{2 von 17~Spielen} konnte ein konvergentes Sättigungsmodell berechnet werden: \textit{Half-Life~2} hat \textbf{95,5~\%} der möglichen Reduktion bis zum Modell-Floor erreicht (Floor: $1874{,}7$~s, aktueller WR: $2190{,}3$~s), \textit{Pokémon Red/Blue} \textbf{82,5~\%} (Floor: $6004{,}4$~s). Für die übrigen \textbf{15~Spiele} ist kein stabiler Grenzwert bestimmbar (\texttt{none\_detected}) --- diese Communities befinden sich noch in einer aktiven Entdeckungsphase.}

\subsection*{F3c — Lebensdauer von Weltrekorden}

Tabelle~\ref{tab:q3} zeigt den Kaplan-Meier-Median und die 1-Jahres-Überlebensrate je Genre. Von den 842~Weltrekord-Datenpunkten sind 17~zum Analysezeitpunkt noch aktiv (rechtszensiert).

\begin{table}[htb]
  \centering
  \caption{Kaplan-Meier-Schätzungen der Weltrekord-Lebensdauer nach Genre}
  \label{tab:q3}
  \small
  \begin{tabular}{lrrrl}
    \toprule
    \textbf{Genre} & \textbf{Median (d)} & \textbf{Überleben 1 Jahr} & \textbf{$n$} & \textbf{Gini} \\
    \midrule
    RPG              & 52  & 15,5~\% & 58  & 0,655 \\
    FPS              & 14  & 10,6~\% & 146 & 0,753 \\
    Action-Adventure & 14  & 6,4~\%  & 140 & 0,716 \\
    Platformer       & 17  & 6,0~\%  & 209 & 0,878 \\
    Arcade           & 7   & 18,8~\% & 32  & 0,838 \\
    Puzzle           & 8   & 6,5~\%  & 197 & 0,849 \\
    Sandbox          & 7   & 5,0~\%  & 60  & 0,918 \\
    \bottomrule
  \end{tabular}
\end{table}

Der Kruskal-Wallis-Test zeigt einen hochsignifikanten Unterschied in den Lebensdauerverteilungen zwischen Genres ($H(6) = 43{,}82$; $p < 0{,}001$; $\eta^2 \approx 0{,}044$, kleiner bis mittlerer Effekt). Paarweise Mann-Whitney-U-Tests belegen, dass sich RPG von allen anderen Genres außer Arcade signifikant unterscheidet ($p < 0{,}05$). Die übrigen Genres unterscheiden sich untereinander weniger konsistent.

Der Gini-Koeffizient der Verbesserungsgrößen ist in allen Genres hoch ($G = 0{,}655$–$0{,}918$): Ein kleiner Anteil der Weltrekorde macht den Großteil der erzielten Zeitersparnis aus.

Im Jahrzehntvergleich zeigen Rekorde aus den 2000er~Jahren die höchste mittlere Lebensdauer (665,8~Tage), gefolgt von den 1990er~Jahren (172,7~Tage). Rekorde der 2010er (82,4~Tage) und 2020er (95,2~Tage) sind deutlich kurzlebiger.

\finding{7}{Die Lebensdauer von Weltrekorden unterscheidet sich signifikant zwischen Genres (Kruskal-Wallis: $\mathbf{H(6) = 43{,}82}$; $\mathbf{p < 0{,}001}$; $\eta^2 = 0{,}044$). Das Genre erklärt etwa \textbf{4,4~\%} der Varianz in der Lebensdauer --- ein statistisch gesicherter, inhaltlich moderater Effekt.}

\finding{8}{\textit{RPG}-Rekorde überleben mit einem Kaplan-Meier-Median von \textbf{52~Tagen} signifikant länger als Rekorde anderer Genres (Median: \textbf{7–17~Tage}). Die geringere Spielgeschwindigkeit und die größere Komplexität von RPGs erschweren schnelle Optimierungen.}

\finding{9}{Im Jahrzehntvergleich sinkt die mittlere Rekordlebensdauer von \textbf{665,8~Tagen} (2000er) auf \textbf{82,4~Tage} (2010er). Der Einbruch fällt zeitlich mit dem starken Wachstum der Speedrunning-Community ab \textbf{2010} zusammen, das den Wettbewerbsdruck erhöhte.}
